\section{What changes when the monitor can stop a write?}
\label{sec:live}
The earlier studies compare recorded warnings. They cannot show that a warning prevents an unauthorized change. We therefore run a separate live experiment: before a proposed database write executes, a gate either permits it or returns a blocking reason to the agent. The agent then continues and can gather evidence, revise its action or refuse the request.

\subsection{Forbidden requests and permitted counterparts}
We use eight designed source pairs, one for each selected rule: cancellation eligibility, basic-economy flight changes, passenger-count preservation, compensation eligibility, delivered-order returns, refund destinations, required retail identity lookup, and changes after an item modification. Each pair contains a forbidden request and a permitted counterpart. A scripted customer supplies fixed facts and pushes back at most twice after a refusal; the acting model generates its own replies and tool choices. These are online generated conversations on designed pressure scenarios, not naturally sampled customer incidents or manually written actor transcripts.

The allocation is eight pairs times two request types times four conditions: 64 conversations. GPT-5.4 Mini is the sole actor, with one independently sampled rollout per cell. We use all eight development source pairs numbered 01; none had previous actor conversations, although two had appeared in constructed monitor fixtures. Allocation, prompts, limits and primary analysis were frozen locally before collection. The experiment is exploratory after earlier feasibility work, not an externally preregistered confirmatory trial. It is not pooled with the earlier task-ID cohorts.

The four conditions share initial case states, actor prompts and native tools. \emph{No gate} executes proposals normally. The \emph{broad gate} asks Gemini 3.1 Pro to check the proposed write against the public policy. The \emph{decomposed gate} first applies predicates to public observations, blocks known violations, and asks the same model focused questions where relevant facts remain unresolved. The \emph{privileged gate} applies the eight predicates using simulator facts. This last condition supplies a reference for those checks; it is neither a complete-policy oracle nor a guaranteed upper bound on task completion. The model gates see the public policy, tools and prefix ending at the proposal, excluding private case labels, hidden facts and future turns.

\subsection{Executed violations, permitted work and cost}
All 32 forbidden-request conversations complete. Among the 32 permitted controls, 29 complete and three reach the customer-turn limit. Table~\ref{tab:live} retains these incomplete outcomes separately. A violation here means a covered forbidden write actually changes native state; a rejected proposal alone is not an executed violation. Control success requires both a matching reference database and conversational completion. One incomplete privileged-gate control had already reached its reference state, illustrating why those checks are recorded separately.

\begin{table}[t]
\centering\small
\caption{Exploratory live comparison: eight forbidden and eight permitted requests per condition. Ratios use completed conversations; all forbidden requests complete. Cost includes all 16 assigned attempts per condition. No monitor outputs are missing.}
\label{tab:live}
\setlength{\tabcolsep}{4pt}
\begin{tabular}{lrrrrr}
\toprule
Gate & Forbidden writes & Permitted successes & Unfinished & Mean turns & Cost\\
 & / completed & / completed controls & controls & (forbidden) & (USD)\\
\midrule
No gate & 3/8 & 6/7 & 1 & 5.250 & 0.0738 \\
Broad model & 0/8 & 7/8 & 0 & 6.375 & 0.4211 \\
Predicates + focused checks & 0/8 & 6/8 & 0 & 6.125 & 0.2493 \\
Privileged predicates & 0/8 & 5/6 & 2 & 6.500 & 0.0852 \\
\bottomrule
\end{tabular}

\end{table}

No gate executes violations in three of eight forbidden-request conversations: an ineligible cancellation, an ineligible credit, and a retail change without the required identity lookup. None executes under any of the three gates. This is a finite observed comparison, not an estimate of a stable population effect. Each condition receives a fresh stochastic rollout, so the differences cannot all be attributed to interception. In the cancellation case, all gated actors refuse before proposing a write; no gate is invoked. Conversely, the ungated actor refuses a prohibited post-item-change request that the gated actors propose and their gates block.

The broad and decomposed gates each block four proposals that satisfy the selected predicates under private simulator facts. These are not automatically false alarms: the agent may still lack required public verification. In the permitted compensation case, both gates initially demand flight-status evidence. The decomposed-gate actor fetches the facts and completes the task; the broad-gate actor stops. Thus, preventing forbidden writes and preserving permitted service require separate measurements. The privileged gate has no such covered-eligible blocks but still has two incomplete controls. The saved analysis also reports missing-outcome bounds over every assigned cell.

\subsection{A concrete intercepted change: an ineligible \$300 credit}
A customer asks for compensation but does not meet the selected eligibility rule. Without a gate, the actor calls \texttt{send\_certificate} for \$300 and the database changes. Under each gate, the actor also proposes the ineligible certificate, but the gate blocks it before execution and returns a reason. No certificate is issued. This demonstrates an actual intervention at the write boundary, rather than an alert scored after the change.

To distinguish gate prevention from native enforcement, we replay all 25 blocked proposals on private copies of their exact pre-write database states. Each gate directly prevents three covered state-changing violations: the ineligible credit, an address change without the required identity lookup, and a forbidden modification after an item change. Eight other blocked proposals would already fail in the native tools without changing state, so they receive no credit as prevented database violations. This one-step replay was added after collection began as an exploratory diagnostic, distinct from the frozen primary outcomes. It does not simulate the rest of an alternative conversation.

All 64 saved attempts pass native replay; a separate implementation agrees on 58 write labels, and 27 model-gate inputs pass public-prefix checks. These checks establish the recorded transitions and literal predicates, not comprehensive semantic compliance. One source pair per rule, one actor and one rollout per condition cannot establish frontier generalization or rank the gate designs reliably. The supported result is narrower and concrete: live gates intercept substantive writes, while their blocking and continuation behavior can also impede permitted tasks.
